% 水星（综述）
% license CCBYSA3
% type Wiki

本文根据 CC-BY-SA 协议转载翻译自维基百科\href{https://en.wikipedia.org/wiki/Mercury_(planet)}{相关文章}。

\begin{figure}[ht]
\centering
\includegraphics[width=8cm]{./figures/8ad510c3701a820c.png}
\caption{} \label{fig_Mercur_1}
\end{figure}
Mercury 是距离太阳最近的行星，也是太阳系中最小的行星。它是一颗岩质行星，具有极其稀薄的大气层，其表面重力略高于火星。水星的表面与地球的月球相似，布满大量撞击坑，并具有由逆断层形成的广阔断崖系统（rupes），以及由抛射物形成的明亮光条结构。其最大撞击盆地——卡洛里斯平原（Caloris Planitia）——直径为 \(1{,}550\,\mathrm{km}\)（\(960\,\mathrm{mi}\)），约为整颗行星直径 \(4{,}880\,\mathrm{km}\)（\(3{,}030\,\mathrm{mi}\)）的三分之一。作为最靠内侧运行的行星，它在地球的天空中总是接近太阳出现，呈现为“晨星”或“昏星”。它也是从地球前往所需 \(\Delta v\) 最大的行星，同时亦是往返太阳系中其他行星时所需能量最高者。

水星的恒星年（\(88.0\) 个地球日）与恒星日（\(58.65\) 个地球日）构成 \(3:2\) 的自转—公转共振。因此，一个太阳日（从日出到下一次日出）约持续 \(176\) 个地球日，即其恒星年的两倍。这意味着水星某一侧将在一个水星年（\(88\) 个地球日）内持续处于日照之下；而在接下来的轨道周期中的另一个 \(88\) 个地球日内，该侧将一直处于黑暗之中，直到下一次日出。行星表面之上存在极度稀薄的外逸层以及微弱但足以偏转太阳风的磁场。由于水星轨道偏心率极高，其表面所接受的日照强度与温度变化极为剧烈：赤道地区夜间温度可达 \(-170^{\circ}\mathrm{C}\)（\(-270^{\circ}\mathrm{F}\)），而白昼可升至 \(420^{\circ}\mathrm{C}\)（\(790^{\circ}\mathrm{F}\)）。由于轴倾角极小，水星两极区域永久处于阴影之中，这强烈暗示其撞击坑内可能存在水冰。

与太阳系其他行星相同，水星形成于约 \(4.5\times10^{9}\) 年前。关于其起源与演化存在诸多相互竞争的假说，其中一些涉及与微行星碰撞及岩石蒸发过程；截至 2020 年代初，水星地质历史的许多宏观细节仍在研究中，或需等待更多深空探测数据。其地幔成分高度均一，这表明水星早期可能存在类似月球的岩浆海。根据当前模型，水星可能具有固态硅酸盐地壳与地幔，其下为固态外核、更深处的液态内核层，以及固态内核。约在七至八十亿年后，随着太阳演化为红巨星，水星预计将被毁灭，金星亦然，地球与月球也可能遭遇同样命运\(^\text{[20]}\)。

水星是一颗“古典行星”，自古以来便被人类持续观测并识别为行星（或“游星”）。英文名称 Mercury 源自古罗马神祇 Mercurius，即商业与沟通之神，也是众神的信使。1974 年，“水手 10 号”（Mariner 10）完成了首次成功飞掠水星的任务，此后 MESSENGER 与 BepiColombo 轨道器也相继对其进行访问与探测。
\subsection{命名}

在历史上，人类会根据水星作为“昏星”或“晨星”出现的不同情况而称呼它不同的名字。到公元前约 350 年，古希腊人已经意识到这两颗“星”实际上是同一个天体\(^\text{[21]}\)。他们称这颗行星为 Στίλβων（Stilbōn），意为“闪烁者”，以及 Ἑρμής（Hermēs），因其迅速而短暂的移动\(^\text{[22]}\)；这一名称在现代希腊语中依然保留（Ερμής Ermis）\(^\text{[23]}\)。罗马人则以他们的信使之神墨丘利（Mercury，拉丁语 Mercurius）为其命名，因为水星在天空中的移动速度比任何其他行星都快\(^\text{[21][24]}\)。不过，根据老普林尼的记载，也有人将这颗行星与阿波罗相联系\(^\text{[25]}\)。

水星的天文学符号是赫尔墨斯权杖（caduceus）的风格化形式；在 16 世纪，该符号加入了一个基督教十字，其形状为：☿\(^\text{[26][27]}\)。
\subsection{物理特性}
\begin{figure}[ht]
\centering
\includegraphics[width=8cm]{./figures/23fb842f59355847.png}
\caption{按比例绘制的水星及内太阳系具有行星质量的天体。从左至右依次为：水星、金星、地球、月球、火星和谷神星。} \label{fig_Mercur_2}
\end{figure}
水星是太阳系中四颗类地行星之一，这意味着它与地球一样，是一颗岩质天体。它是太阳系中最小的行星，赤道半径为 \(2{,}439.7\,\mathrm{km}\)（\(1{,}516.0\,\mathrm{mi}\)）\(^\text{[4]}\)。水星的体积也小于——但质量大于——太阳系中最大的两个天然卫星：木卫三（盖尼米得）和土卫六（泰坦）。水星的组成大约为 70\% 的金属物质与 30\% 的硅酸盐物质\(^\text{[28]}\)。
\subsubsection{内部结构}
\begin{figure}[ht]
\centering
\includegraphics[width=8cm]{./figures/697a86a5c3dcc15c.png}
\caption{水星的内部结构与磁场} \label{fig_Mercur_3}
\end{figure}
水星似乎具有一层固态硅酸盐地壳与地幔，它们覆盖着固态金属外核层、更深的液态内核层，以及固态内核\(^\text{[29][30]}\)。富铁内核的成分仍不确定，但很可能包含镍、硅，或许还有硫和碳，并含有微量其他元素\(^\text{[31]}\)。水星的密度是太阳系中第二高的，为 \(5.427\,\mathrm{g/cm^{3}}\)，仅略低于地球的 \(5.515\,\mathrm{g/cm^{3}}\)\(^\text{[4]}\)。如果将两颗行星的重力压缩效应排除，水星的构成物质将比地球更致密：其未压缩密度约为\(5.3\,\mathrm{g/cm^{3}}\)，而地球为 \(4.4\,\mathrm{g/cm^{3}}\)\(^\text{[32]}\)。水星的密度可用于推断其内部结构。尽管地球的高密度在很大程度上源于重力压缩（尤其在核心处），但水星要小得多，其内部区域并未受到同等程度的压缩。因此，为了呈现如此高的密度，其核心必须既巨大又富含铁\(^\text{[33]}\)。

基于满足惯性矩因子 \(0.346\pm0.014\) 的内部结构模型，水星的核心半径估计为 \(2{,}020\pm30\,\mathrm{km}\)（\(1{,}255\pm19\,\mathrm{mi}\)）\(^\text{[9][34]}\)。因此，水星核心占其总体积约 57\%；相比之下，地球仅为 17\%。2007 年发表的研究表明水星具有熔融内核\(^\text{[35][36]}\)。地幔—地壳层总厚度约为 \(420\,\mathrm{km}\)（\(260\,\mathrm{mi}\)）\(^\text{[37]}\)。关于地壳厚度，各研究推算存在差异：Mariner~10 与 MESSENGER 探测数据给出的估计为 \(35\,\mathrm{km}\)（\(22\,\mathrm{mi}\)），而基于 Airy 均衡假说的模型则给出 \(26\pm11\,\mathrm{km}\)（\(16.2\pm6.8\,\mathrm{mi}\)）\(^\text{[38][39][40]}\)。水星表面的一个显著特征是大量狭长的脊状地形，可延伸数百公里。一般认为这些构造形成于水星核心与地幔冷却并收缩的时期，而那时地壳已处于固态\(^\text{[41][42][43]}\)。

水星的核心含铁量在太阳系行星中最高，为解释这一特征已提出多种理论。最广为接受的观点认为，水星最初的金属—硅酸盐比例与普通球粒陨石相似，而后者被认为代表太阳系岩质物质的典型组成；当时水星的质量大约为其当前质量的 2.25 倍\(^\text{[44]}\)。在太阳系早期，水星可能被一颗质量约为其六分之一、直径达数千公里的微行星撞击\(^\text{[44]}\)。这一撞击会剥离其大部分原始地壳与地幔，使核心成为主要残留成分\(^\text{[44]}\)。一种类似的机制，即巨撞假说（giant impact hypothesis），也被用于解释地球月球的形成\(^\text{[44]}\)。

另一种假说认为，水星可能形成于太阳星云的早期阶段，那时太阳的能量输出尚未稳定。其初始质量可能是当前质量的两倍，但随着原太阳（protosun）收缩，水星附近的温度可能达到 \(2{,}500\sim3{,}500\,\mathrm{K}\)，甚至高达 \(10{,}000\,\mathrm{K}\)\(^\text{[45]}\)。在如此温度下，水星的大部分表面岩石会被汽化，形成“岩石蒸汽”大气，并可能被太阳风带走\(^\text{[45]}\)。第三种假说则认为，太阳星云对水星聚积的固体颗粒施加阻力，导致较轻的颗粒从聚积物质中被移除，从而未被水星捕获\(^\text{[46]}\)。

每一种假说均预测不同的表面化学组成，因此已有两项深空任务专门针对这一问题进行观测。第一项任务 MESSENGER（2015 年结束）发现水星表面的钾与硫含量高于预期，这表明巨撞假说与地壳—地幔蒸发模型可能并未发生，因为这些事件的极端高温会使钾与硫完全挥发\(^\text{[47]}\)。BepiColombo 将于 2025 年抵达水星，并将继续对这些假说进行检验\(^\text{[48]}\)。迄今为止的结果似乎支持第三种假说；然而，仍需进一步分析数据\(^\text{[49]}\)。
\subsubsection{表面地质}
\begin{figure}[ht]
\centering
\includegraphics[width=8cm]{./figures/02332822930c85ad.png}
\caption{利用 NASA 和 USGS 数据在 Blender 中生成的水星渲染图} \label{fig_Mercur_4}
\end{figure}
水星的表面在外观上与月球相似，呈现广阔的类月海平原以及大量撞击坑，这表明其地质活动已停滞达数十亿年之久。与火星或月球相比，水星表面更加不均质；火星与月球都具有大面积的相似地质单元，例如月海与高原\(^\text{[50]}\)。反照率特征是指具有显著不同反射率的区域，其中包括撞击坑、由此形成的抛射物以及光条系统。较大的反照率特征通常对应反射率更高的平原区域\(^\text{[51]}\)。水星具有“皱褶脊”（dorsa）、类似月球的高地、山脉（montes）、平原（planitiae）、断崖（rupes）以及峡谷（valles）\(^\text{[52][53]}\)。
\begin{figure}[ht]
\centering
\includegraphics[width=8cm]{./figures/8c611b128a504de7.png}
\caption{MESSENGER 对水星表面进行的 MASCS 光谱扫描} \label{fig_Mercur_5}
\end{figure}
水星的地幔在化学成分上具有不均一性，这表明该行星在早期历史中经历过岩浆海阶段。矿物的结晶与对流翻转导致形成了一层具有化学分层、成分不均的地壳，而这种大尺度的化学组成差异也可在其表面观测到。地壳含铁量低但含硫量高，这源于水星早期比其他类地行星更强的化学还原环境。水星表面主要由贫铁的辉石与橄榄石组成，分别以顽辉石（enstatite）与镁橄榄石（forsterite）为代表，同时还包含富钠的斜长石，以及由镁、钙与硫化铁混合构成的矿物。地壳中反照率较低的区域含有高含量的碳，最可能以石墨的形式存在\(^\text{[54][55]}\)。

水星地表特征的命名来自多种来源，并依据 IAU 行星命名系统设定。以人物命名的仅限于已故者。撞击坑以在艺术、音乐、绘画与文学等领域做出杰出或基础性贡献的人物命名。脊或皱褶（dorsa）以对水星研究有贡献的科学家命名。洼地或槽谷（fossae）以建筑作品命名。山脉（montes）以各语言中表示“热”的词语命名。平原（planitiae）以各语言中“水星”的名称命名。断崖（rupēs）以科学考察的船只命名。峡谷（valles）以古代被废弃的城市、城镇或聚落命名\(^\text{[56]}\)。

\textbf{撞击盆地与撞击坑}
\begin{figure}[ht]
\centering
\includegraphics[width=8cm]{./figures/cc1a7ee3ff642536.png}
\caption{卡洛里斯盆地附近火山平原（橙色）中曼奇陨石坑（左）、桑德陨石坑（中）与坡陨石坑（右）的增强彩色影像} \label{fig_Mercur_6}
\end{figure}
在大约 46 亿年前水星形成期间及形成后不久，它曾遭受彗星和小行星的猛烈轰击，并可能在随后一个独立阶段（称为晚期重轰炸时期）再次经历强烈撞击，该时期约在 38 亿年前结束\(^\text{[57]}\)。在这一剧烈的撞击坑形成时期，水星的整个表面都遭到撞击\(^\text{[53]}\)，而不存在任何能够减速撞击体的大气层更进一步加剧了这一过程\(^\text{[58]}\)。在此阶段，水星曾经历火山活动；大量岩浆填充盆地，形成类似月球月海的平滑平原\(^\text{[59][60]}\)。其中一个最为特殊的撞击坑是阿波罗多鲁斯（Apollodorus），又称“蜘蛛（the Spider）”，其内部拥有一系列从撞击中心向外辐射的槽谷\(^\text{[61]}\)。

水星上的撞击坑直径范围极广，从小型碗状洼地到横跨数百公里的多环撞击盆地皆有。它们呈现出不同程度的风化状态，从相对年轻、带有光条的撞击坑，到高度侵蚀、几乎只剩残迹的古老撞击坑。与月球撞击坑相比，水星撞击坑的一个细微差异在于其抛射物覆盖区域更小，这是因为水星表面重力更强所致\(^\text{[62]}\)。根据国际天文学联合会（IAU）的命名规则，每一个新撞击坑必须以在世至少五十年成名、并且在命名前至少已逝世三年的艺术家命名\(^\text{[63]}\)。
\begin{figure}[ht]
\centering
\includegraphics[width=8cm]{./figures/61b892c36d341211.png}
\caption{卡洛里斯盆地的俯视图/卡洛里斯盆地的透视图——高处（红色）；低处（蓝色）} \label{fig_Mercur_7}
\end{figure}
已知最大的撞击坑是卡洛里斯平原（Caloris Planitia），又称卡洛里斯盆地（Caloris Basin），其直径为 \(1{,}550\,\mathrm{km}\)（\(960\,\mathrm{mi}\)）\(^\text{[64]}\)。形成卡洛里斯盆地的撞击极其强大，以至于引发了岩浆喷发，并在撞击坑周围形成了一圈约 \(2\,\mathrm{km}\)（\(1.2\,\mathrm{mi}\)）高的同心山脉。卡洛里斯盆地的底部被一种地质上独特的平坦平原所填充，其中由脊与裂隙构成近似多边形的分布格局。目前尚不清楚这些构造是撞击诱发的火山熔岩流所形成，还是由撞击熔融物的大型熔融层所造成\(^\text{[62]}\)。

在卡洛里斯盆地的对跖点位置有一大片奇特的丘陵地形，被称为“奇异地形（Weird Terrain）”。其中一个假说认为，该地形源自卡洛里斯撞击产生的冲击波沿水星传播，并在对跖点（相隔 180 度）汇聚，从而导致强烈的应力使地表出现破碎\(^\text{[65]}\)。另一种观点则认为，该地形可能由撞击抛射物在盆地对跖点汇聚堆积而形成\(^\text{[66]}\)。
\begin{figure}[ht]
\centering
\includegraphics[width=8cm]{./figures/03d5b8f93bb45933.png}
\caption{在这幅水星边缘的图像中，托尔斯泰盆地位于底部。} \label{fig_Mercur_8}
\end{figure}
总体而言，水星已确认有 46 个撞击盆地\(^\text{[67]}\)。其中一个显著的盆地是托尔斯泰盆地（Tolstoj Basin），其宽度为 \(400\,\mathrm{km}\)（\(250\,\mathrm{mi}\)），为多环结构，其抛射物覆盖层可从盆地边缘延伸至 \(500\,\mathrm{km}\)（\(310\,\mathrm{mi}\)）之外，盆地底部则被平滑平原物质所填充。贝多芬盆地（Beethoven Basin）具有与托尔斯泰盆地类似规模的抛射物覆盖层，其盆地边缘直径为 \(625\,\mathrm{km}\)（\(388\,\mathrm{mi}\)）\(^\text{[62]}\)。与月球相似，水星表面也可能长期受到太空风化过程的影响，包括太阳风与微流星体撞击\(^\text{[68]}\)。

\textbf{平原}

水星上存在两类地质上截然不同的平原区域\(^\text{[62][69]}\)。位于撞击坑之间的缓坡起伏丘陵平原是水星最古老的可见地表，它们早于高度密集撞击坑区域的形成\(^\text{[62]}\)。这些坑间平原似乎掩盖了许多更早期的撞击坑，并且在小于约 \(30\,\mathrm{km}\)（\(19\,\mathrm{mi}\)）的撞击坑数量上普遍偏少\(^\text{[69]}\)。

光滑平原是广布的平坦区域，它们填充了大小不一的洼地，外观上与月球月海高度相似。但不同于月球月海，水星的光滑平原与更古老的坑间平原具有相同的反照率。尽管缺乏明确无误的火山特征，但这些平原的分布位置以及圆滑、舌状的外形都强烈支持其火山起源\(^\text{[62]}\)。所有光滑平原都形成于卡洛里斯盆地之后，这一点由其撞击坑密度明显低于卡洛里斯盆地抛射物覆盖层所证明\,\(^\text{[62]}\)。

\textbf{压缩构造特征}

水星表面最独特的特征之一是大量的压缩褶皱或断崖（rupes），它们在平原上纵横交错。这类构造在月球上也存在，但在水星上更为显著\(^\text{[70]}\)。随着水星内部冷却，其体积收缩，导致地表开始变形，形成皱褶脊与与逆断层相关的舌状断崖。这些断崖的长度可达 \(1{,}000\,\mathrm{km}\)（\(620\,\mathrm{mi}\)），高度可达 \(3\,\mathrm{km}\)（\(1.9\,\mathrm{mi}\)）\(^\text{[71]}\)。这些压缩构造叠加在其他地貌之上，例如撞击坑和光滑平原，说明它们形成得更为晚近\(^\text{[72]}\)。对这些构造的制图表明，水星半径的总体收缩量大约在 \(1\sim7\,\mathrm{km}\)（\(0.62\sim4.35\,\mathrm{mi}\)）之间\(^\text{[73]}\)。主要逆断层系统的大部分活动可能在约 36–37 亿年前结束\(^\text{[74]}\)。已经发现了小尺度的逆断层断崖，其高度仅数十米，长度为数公里，它们看起来不到 5,000 万年历史，说明内部压缩与由此引发的地质活动一直持续至今\(^\text{[71][73]}\)。

\textbf{火山活动}
\begin{figure}[ht]
\centering
\includegraphics[width=8cm]{./figures/74c043b705258e85.png}
\caption{毕加索陨石坑东侧的弧形洼坑可能由岩浆管塌陷形成。} \label{fig_Mercur_9}
\end{figure}
有证据表明，水星上存在来自低矮盾状火山的火山碎屑流\(^\text{[75][76][77]}\)。目前已识别出 51 处火山碎屑沉积物\(^\text{[78]}\)，其中 90\% 分布在撞击坑内部\(^\text{[78]}\)。对这些承载火山碎屑沉积物的撞击坑风化状态的研究表明，水星上的火山碎屑活动在相当长的一段时间内持续发生\(^\text{[78]}\)。

卡洛里斯盆地西南边缘内部的一个“无缘洼地”（rimless depression）由至少九个重叠的火山喷口组成，每个喷口的直径可达 \(8\,\mathrm{km}\)（\(5.0\,\mathrm{mi}\)）。因此，这里是一座“复合火山（compound volcano）”\(^\text{[79]}\)。喷口底部至少比其边缘低 \(1\,\mathrm{km}\)（\(0.62\,\mathrm{mi}\)），其形态更类似于由爆炸式喷发雕刻而成的火山口，或由于岩浆回撤至通道深处形成空洞而发生坍塌后所形成的构造\(^\text{[79]}\)。科学家虽然无法精确定年该火山复合体的年龄，但推测其可能达到十亿年量级\(^\text{[79]}\)。
\subsubsection{表面环境与外逸层}
\begin{figure}[ht]
\centering
\includegraphics[width=8cm]{./figures/1b014d9a9456b7ab.png}
\caption{水星北极的合成图像，其中大量水冰储存在永久阴影的陨石坑内\(^\text{[80]}\)。} \label{fig_Mercur_10}
\end{figure}
水星的表面温度范围为 \(100\sim700\,\mathrm{K}\)（\(-173\sim427\,^{\circ}\mathrm{C}\)；\(-280\sim800\,^{\circ}\mathrm{F}\)）\,[81]。由于缺乏大气层以及赤道与两极之间的陡峭温度梯度，其两极温度从未超过 \(180\,\mathrm{K}\)\(^\text{[15]}\)。在近日点时，赤道的太阳直射点位于西经 \(0^{\circ}\) 或 \(180^{\circ}\)，温度可升至约 \(700\,\mathrm{K}\)。在远日点时，太阳直射点位于西经 \(90^{\circ}\) 或 \(270^{\circ}\)，温度仅达到约 \(550\,\mathrm{K}\)\(^\text{[82]}\)。在水星暗侧，平均温度约为 \(110\,\mathrm{K}\)\(^\text{[15][83]}\)。水星表面的日照强度介于 4.59 至 10.61 倍太阳常数之间（\(1{,}370\,\mathrm{W\,m^{-2}}\)）\(^\text{[84]}\)。

尽管水星表面的白昼温度普遍极高，但观测结果强烈表明水星上存在冰（冻结水）。两极深坑的坑底从未受到直接阳光照射，温度始终保持在 \(102\,\mathrm{K}\) 以下，远低于全球平均水平\(^\text{[85]}\)。这形成了冰得以积累的“冷阱”。水冰具有很强的雷达反射性，1990 年代初，使用 70 米口径的戈德斯通太阳系雷达（Goldstone Solar System Radar）与 VLA 的观测显示，两极附近存在高雷达反射区域\(^\text{[86]}\)。虽然这些反射区域并非只有冰才能造成，但天文学家认为冰是最可能的解释\(^\text{[87]}\)。随后，通过 MESSENGER 对北极陨石坑的成像，水冰的存在得到确认\(^\text{[80]}\)。

这些含冰陨石坑区域估计包含约 \(10^{14}\sim10^{15}\,\mathrm{kg}\) 的冰\(^\text{[88]}\)，并可能被一层抑制升华的风化层覆盖\(^\text{[89]}\)。相比之下，地球南极冰盖的质量约为 \(4\times10^{18}\,\mathrm{kg}\)，而火星南极冠则包含约 \(10^{16}\,\mathrm{kg}\) 的水\(^\text{[88]}\)。水星表面冰的来源尚不明确，但最可能的两种机制为：来自水星内部的气体释放（outgassing）以及彗星撞击所带来的沉积\(^\text{[88]}\)。

水星的体积太小、温度太高，其引力无法在长时间尺度上留住任何显著的大气层；但它确实具有一种由地表约束的稀薄外逸层\(^\text{[90]}\)，其地表压力低于约 \(0.5\,\mathrm{nPa}\)（\(0.005\,\mathrm{picobar}\)）\(^\text{[4]}\)。该外逸层包含氢、氦、氧、钠、钙、钾、镁、硅、氢氧根等多种成分\(^\text{[18][19]}\)。这种外逸层并不稳定——原子不断从多种来源中损失又被重新补充。氢原子与氦原子可能来自太阳风，它们首先扩散进入水星的磁层，然后再逃逸回太空。水星地壳中元素的放射性衰变是氦以及钠与钾的另一来源。水蒸气也存在，其来源包括彗星撞击地表、太阳风中的氢与岩石中的氧通过溅射（sputtering）过程生成水，以及永久阴影极区陨石坑中的水冰升华等多种机制。检测到高含量的水相关离子如 \(\mathrm{O^{+}}\)、\(\mathrm{OH^{-}}\) 与 \(\mathrm{H_{3}O^{+}}\) 是一个意外发现\(^\text{[91][92]}\)。由于这些离子在水星周围空间环境中的含量很高，科学家推测它们是被太阳风从表面或外逸层中击出\(^\text{[93][94]}\)。

钠、钾与钙在 1980–1990 年代被发现存在于水星外逸层中，一般认为主要来源于微流星体撞击地表岩石所造成的蒸发\(^\text{[95]}\)，包括目前来自恩克彗星（Comet Encke）的尘埃\(^\text{[96]}\)。2008 年，MESSENGER 探测器又发现了镁\(^\text{[97]}\)。研究显示，在某些时刻，钠的发射增强出现在与水星磁极对应的位置，这表明磁层与行星表面存在相互作用\(^\text{[98]}\)。

根据 NASA 的说法，水星并不适合地球类生命生存。它具有地表边界外逸层而非分层大气，同时表面温度极端且太阳辐射强烈，不太可能有任何生命能够在这种环境下存活\(^\text{[99]}\)。然而，水星地下部分的某些区域可能曾经具有适宜生存的条件，进而可能孕育过生命形式，尽管此类生命很可能只是原始微生物\(^\text{[100][101][102]}\)。
\subsubsection{磁场与磁层}
\begin{figure}[ht]
\centering
\includegraphics[width=8cm]{./figures/6666002d08218c01.png}
\caption{显示水星磁场相对强度的图示} \label{fig_Mercur_11}
\end{figure}
尽管体积小且自转缓慢（自转周期为 59 天），水星却拥有显著且似乎是全球性的磁场。根据水手 10 号（Mariner~10）的测量，其磁场强度约为地球的 1.1\%。水星赤道处的磁场强度约为 \(300\,\mathrm{nT}\)\(^\text{[103][104]}\)。与地球相似，水星的磁场呈偶极结构\(^\text{[98]}\)，且与行星自转轴近乎对齐（偶极倾角为 \(10^{\circ}\)，相比之下地球为 \(11^{\circ}\)）\(^\text{[105]}\)。来自水手 10 号与 MESSENGER 两个探测器的测量均表明，这一磁场的强度与形状保持稳定\(^\text{[105]}\)。

该磁场很可能由发电机效应（dynamo effect）产生，与地球磁场的形成方式类似\(^\text{[35][106]}\)。这一发电机效应应当源于水星富铁液态内核的循环流动。水星高轨道偏心率所造成的强烈潮汐加热效应，有助于维持内核部分区域保持液态，这是发电机效应得以运作的必要条件\(^\text{[107][108]}\)。

水星的磁场足以使太阳风绕行其外部，从而形成磁层。尽管水星的磁层小到可以完全容纳在地球磁层范围内\(^\text{[98]}\)，但其强度足以捕获太阳风等离子体。这种过程会促进水星表面的太空风化作用\(^\text{[105]}\)。水手 10 号在探测中发现，水星夜侧磁层中存在低能等离子体。其磁尾中出现的高能粒子爆发则表明水星的磁层具有动态特性\(^\text{[98]}\)。

在 2008 年 10 月 6 日进行的第二次飞掠中，MESSENGER 发现水星的磁场可能极为“漏散”（leaky）。探测器遇到直径可达 \(800\,\mathrm{km}\)（约为行星半径的三分之一）的磁场“龙卷风”——即扭曲的磁力线束，它们将水星磁场直接连接至行星际空间。这些扭曲磁通管（flux transfer events，简称 FTE）在水星磁场屏障中形成开放窗口，使太阳风能够通过磁重联进入并直接冲击水星表面\(^\text{[109]}\)。类似现象也发生在地球磁场中。MESSENGER 的观测显示，水星的磁重联速率是地球的十倍，而水星更靠近太阳这一因素仅能解释该速率约三分之一的增强\(^\text{[109]}\)。
\subsection{轨道、自转与经度}
水星在太阳系所有行星中具有最偏心的轨道；其轨道偏心率为 0.21，与太阳的距离在 \(46{,}000{,}000\sim70{,}000{,}000\,\mathrm{km}\)（\(29{,}000{,}000\sim43{,}000{,}000\,\mathrm{mi}\)）之间变化。水星完成一次公转需 87.969 个地球日。图示展示了轨道偏心率的影响，将水星轨道与具有相同半长轴的圆形轨道叠加进行比较。当接近近日点时，水星的速度明显更高，从其在每个 5 天区间内覆盖的更大距离即可看出。在图中，水星距太阳的变化通过行星大小的变化表示，行星大小与其距太阳的距离成反比。

水星与太阳距离的变化使其表面受到潮汐隆起的影响，这些潮汐隆起由太阳产生，其强度约为地球上月球潮汐的 17 倍\(^\text{[110]}\)。结合水星绕自转轴的 \(3:2\) 自转—公转共振，这也导致其表面温度呈现复杂变化\(^\text{[28]}\)。这一共振使得水星的一个太阳日（两次太阳过中天之间的时间）恰好等于两个水星年，即约等于 176 个地球日\(^\text{[111]}\)。

水星轨道相对于地球轨道平面（黄道面）的倾角为 \(7^{\circ}\)，这是太阳系八大行星中最大的倾角\(^\text{[112]}\)。因此，水星凌日现象只有在其位于地球与太阳之间、并同时通过黄道面时才会发生，即发生在每年的五月或十一月。平均而言，这一现象大约每七年出现一次\(^\text{[113]}\)。

水星的自转轴倾角几乎为零\(^\text{[114]}\)，目前最佳测量值低至 \(0.027^{\circ}\)\(^\text{[115]}\)。这个数值显著小于轴倾角排名第二小的木星（\(3.1^{\circ}\)）。因此，对于位于水星两极的观察者而言，太阳中心从未升至距离地平线超过 2.1 角分的高度\,\(^\text{[115]}\)。相比之下，从水星上看到的太阳角直径范围从 \(1\tfrac{1}{4}\) 度到 2 度\(^\text{[116]}\)。

在水星表面的某些地点，观察者能够看到太阳升起到略高于地平线三分之二的位置，然后倒退落下，随后再次升起，而这一切都发生在同一个水星日之内\(^\text{[a]}\)。其原因在于：在近日点前约四个地球日，水星的角轨道速度与其角自转速度相等，使太阳的视运动停止；更接近近日点时，水星的角轨道速度超过角自转速度。因此，对于假想的水星表面观察者来说，太阳似乎会呈现逆行运动。近日点之后四个地球日，太阳正常的视运动恢复\(^\text{[28]}\)。如果水星曾处于同步自转状态，将会出现类似效应：在一个公转周期内，自转速度的交替增减会产生经度方向 \(23.65^{\circ}\) 的天平动现象\(^\text{[117]}\)。

基于同样原因，水星赤道上相隔 180 度的两个地理点，在交替的水星年（即一个水星日）中的近日点附近，会出现如下现象：太阳从头顶经过，然后逆转其视运动，再次从头顶经过；之后太阳再一次逆转，第三次从头顶经过。整个过程持续约 16 个地球日。在另一个交替的水星年，同样的现象会发生在另一对点上。逆行运动的幅度很小，因此整体效果是：在一段两到三周的时间里，太阳几乎在头顶附近静止不动，并且由于水星位于近日点（最接近太阳），太阳亮度达到最强。由于长时间暴露在最强烈的阳光下，这两个地点成为水星表面最热的地方。最高温度出现在太阳超过正午约 \(25^{\circ}\) 之后，这是由于昼夜温度滞后效应造成的，即在日出后约 0.4 个水星日与 0.8 个水星年处达到峰值\(^\text{[118]}\)。相反，赤道上另有两个地点，与前述地点相隔 90 度经度，在交替的年份里，只有当水星位于远日点时太阳才会从头顶经过，此时太阳在水星天空中的视运动相对快速。这些地点正是前一段所述太阳在越过地平线时出现逆行视运动的位置，因此它们接收到的太阳热量远小于前述近日点附近的那两个地点\(^\text{[119]}\)。

水星平均每 116 个地球日达到一次下合（与地球最近）\(^\text{[4]}\)，但由于轨道偏心率高，这一周期可在 105 至 129 天之间变化。水星与地球的最近距离可达 \(82{,}200{,}000\,\mathrm{km}\)（0.549\,AU；5110 万英里），且这一最近距离正缓慢缩小：下次达到 \(82{,}100{,}000\,\mathrm{km}\)（5100 万英里）将在公元 2679 年，其后在公元 4487 年将达到 \(82{,}000{,}000\,\mathrm{km}\)（5100 万英里），但直到公元 28,622 年前，水星都不会比 8000 万公里（5000 万英里）更接近地球\(^\text{[120]}\)。从地球观测，水星的逆行期（出现在下合前后）可在 8 至 15 天之间变化，而这一大范围的变动也源于其高轨道偏心率\(^\text{[28]}\)。总体而言，由于水星是距离太阳最近的行星，从长期平均来看，它最常是距离地球最近的行星\(^\text{[121][122]}\)，并且在同一意义下，它也是太阳系中所有其他行星最常距离它们最近的行星\(^\text{[121][123][124][b]}\)。
\subsubsection{经度约定}
水星的经度约定将经度零度定义在其表面两个最热地点之一，如前文所述。然而，当水手 10 号（Mariner~10）首次访问该区域时，该零度子午线区域处于黑暗之中，因此无法选择具体地表特征来精确定义子午线的位置。于是，选择了其更西侧的一处小型陨石坑——Hun~Kal，作为经度测量的精确参考点\(^\text{[125][126]}\)。Hun~Kal 的中心定义了西经 \(20^{\circ}\) 的子午线。1970 年国际天文学联合会（IAU）的决议建议，水星的经度应按向西方向为正进行测量\(^\text{[127]}\)。因此，赤道上两个最热的地点位于西经 \(0^{\circ}\) 与 \(180^{\circ}\)，而赤道上最冷的地点位于西经 \(90^{\circ}\) 与 \(270^{\circ}\)。然而，MESSENGER 探测计划采用的是东向为正的经度约定\(^\text{[128]}\)。
\subsubsection{自旋—轨道共振}
\begin{figure}[ht]
\centering
\includegraphics[width=8cm]{./figures/e889dfbdf76e1612.png}
\caption{在完成一次公转后，水星自转了 1.5 圈，因此在完成两次完整公转后，同一半球将再次被太阳照亮。} \label{fig_Mercur_12}
\end{figure}
多年来，人们一直认为水星与太阳存在同步潮汐锁定，即水星每公转一周便自转一周，因此始终以同一面对着太阳，就像月球的同一面始终朝向地球一样。1965 年的雷达观测证明，这一观点是不正确的：水星实际上处于 \(3:2\) 自旋—轨道共振状态，即每绕太阳公转两周，自转三周。水星轨道的偏心率使这一共振得以稳定——在近日点处，太阳潮汐力最强，太阳在水星天空中几乎保持静止\(^\text{[129]}\)。

这种 \(3:2\) 共振潮汐锁定，由沿着水星偏心轨道变化的潮汐力作用于水星质量分布中的永久偶极分量而稳定\(^\text{[130]}\)。在圆形轨道中并不存在这种潮汐力变化，因此唯一能稳定的共振是 \(1:1\)（例如地月系统），此时潮汐力沿“中心天体方向”拉伸天体，施加的力矩会使天体的最小惯性轴（“最长”轴，即对应前述偶极方向的轴）始终指向中心天体。然而，若轨道偏心率显著，如同水星轨道，潮汐力在近日点达到最大，因此能够稳定 \(3:2\) 这样的共振，使得行星在经过近日点时，其最小惯性轴大致指向太阳\(^\text{[130]}\)。

天文学家最初误以为水星是同步自转的原因在于：每当水星处于最适宜观测的位置时，它恰好总是位于其 \(3:2\) 共振的同一相位，因此表现为总是朝向同一面。这是因为巧合地，水星的自转周期几乎正好是其相对于地球的会合周期的一半。由于 \(3:2\) 自旋—轨道共振，一个太阳日约等于 176 个地球日\(^\text{[28]}\)；一个恒星日（自转周期）约为 58.7 个地球日\(^\text{[28]}\)。

数值模拟显示，由于来自其他行星的摄动，水星轨道偏心率在数百万年尺度上呈现混沌变化，可从近乎零（圆轨道）增加到超过 0.45\(^\text{[28][131]}\)。人们曾认为，这可解释水星为何落入 \(3:2\) 共振（而不是更常见的 \(1:1\) 共振），因为高偏心率时期更容易导致这种状态的捕获\(^\text{[132]}\)。然而，基于实际潮汐响应的精确模拟表明，水星在其历史的非常早期——形成后约 2,000 万年（更可能是 1,000 万年）内——便已被捕获到 \(3:2\) 共振状态\(^\text{[133]}\)。

数值模拟还显示，在未来，由于水星轨道近日点与木星近日点的长期轨道共振相互作用，其轨道偏心率可能增加到足以使轨道失稳的程度，在未来五十亿年内大约有 1\% 的概率发生这种情况。如果发生这种不稳定，水星可能坠入太阳、与金星相撞、被太阳系抛射出去，甚至扰乱整个内太阳系\(^\text{[134][135]}\)。
\subsubsection{近日点进动}
\begin{figure}[ht]
\centering
\includegraphics[width=8cm]{./figures/a9bd2f7841fd09f8.png}
\caption{水星轨道的近心点进动} \label{fig_Mercur_13}
\end{figure}
1859 年，法国数学家兼天文学家乌尔班·勒维耶（Urbain Le~Verrier）报告称，水星绕太阳轨道的缓慢进动无法完全用牛顿力学及已知行星的摄动来解释。他提出的一种可能解释是：在水星轨道以内可能存在另一颗行星（或一系列较小的“微粒”），它们的引力可解释这种摄动\(^\text{[136]}\)。其他被考虑的解释包括太阳具有轻微的扁率。由于天文学家通过天王星轨道摄动成功发现海王星，人们对这一类推非常信任，并将这颗假想行星命名为伏尔坎（Vulcan），但始终未被发现\(^\text{[137]}\)。

观测到的水星近日点进动相对于地球为每世纪 \(5{,}600\) 角秒（\(1.5556^{\circ}\)），或相对于惯性参考系 ICRF 为 \(574.10\pm0.65\) 角秒每世纪\(^\text{[138]}\)。牛顿力学在考虑所有行星的摄动影响并包含太阳扁率所致的每世纪 \(0.0254\) 角秒的贡献后，预测水星近日点进动相对于地球为 \(5{,}557\) 角秒（\(1.5436^{\circ}\)）每世纪，或相对于 ICRF 为 \(531.63\pm0.69\) 角秒每世纪\(^\text{[138]}\)。20 世纪初，阿尔伯特·爱因斯坦的广义相对论给出了观测进动的理论解释，即将引力形式化为由时空曲率所介导的现象。该效应非常小：对水星而言仅为每世纪 \(42.980\pm0.001\) 角秒（或每年 0.43 角秒，或每轨道周期 0.1035 角秒）。因此，水星需要略多于 1,250 万次轨道运行（约 300 万年）才能累积形成一次完整的额外超前转动。类似的，但更小的效应也存在于太阳系其他天体上：金星为每世纪 8.6247 角秒，地球为 3.8387 角秒，火星为 1.351 角秒，小行星 1566~Icarus 为 10.05 角秒\(^\text{[139][140]}\)。
\subsection{观测}
\begin{figure}[ht]
\centering
\includegraphics[width=6cm]{./figures/a8855b4cf009773a.png}
\caption{水手 10 号于 1974 年拍摄的图像拼接} \label{fig_Mercur_14}
\end{figure}
水星的视星等计算表明，在上合日附近可达到 \(-2.48\)（比天狼星更亮），而在下合日附近则可降至 \(+7.25\)（低于裸眼可见极限）\(^\text{[16]}\)。其平均视星等为 0.23，而标准差为 1.78，是所有行星中最大的。水星在上合日的平均视星等为 \(-1.89\)，而在下合日为 \(+5.93\)\(^\text{[16]}\)。由于水星非常接近太阳，其观测常被太阳的强光所掩盖，因此观测非常困难。水星只能在晨昏薄明的短暂时间内被观测到\(^\text{[141]}\)。

从地面望远镜观测水星时，所能看到的仅是一个亮度有限的部分照亮的盘面。哈勃太空望远镜完全无法观测水星，因为其安全程序禁止望远镜指向离太阳过近的位置\(^\text{[142]}\)。由于地球在一个水星年中相对转动了 0.15 圈，这意味着一个七水星年的周期（\(0.15 \times 7 \approx 1.0\)）；因此在第七个水星年中，水星几乎精确（提前 7 天）重复七水星年前所呈现的现象序列\(^\text{[143]}\)。

与月球和金星类似，水星从地球上看也呈现相位变化。在下合日时呈“新相”，在上合日时呈“满相”。在这两种情况下，由于被太阳遮蔽，水星从地球上通常不可见\(^\text{[141]}\)，除非在新相期间发生凌日现象。严格来说，当水星处于满相时，从地球看最为明亮。尽管水星在满相时距离地球最远，但由于可见照明面积更大，加上冲增强效应，其亮度仍超过距离较近时的反照亮度\(^\text{[144]}\)。金星则相反：当它呈现弦月形时最亮，因为在此相位下它距离地球比盈相更加接近\(^\text{[144][145]}\)。
\begin{figure}[ht]
\centering
\includegraphics[width=8cm]{./figures/51ad7702caf71357.png}
\caption{显示北极区域最高温度的伪彩色地图} \label{fig_Mercur_15}
\end{figure}
\begin{figure}[ht]
\centering
\includegraphics[width=8cm]{./figures/94480d99e744851f.png}
\caption{从加利福尼亚州圣何塞观测到的水星（左下），同框可见金星与月球。} \label{fig_Mercur_16}
\end{figure}
水星在上弦和下弦位置时最适合观测，尽管这些相位的亮度较弱。上弦相与下弦相分别出现在其相对于太阳的最大东、最大西伸角时。在这两个时刻，水星与太阳的角距从近日点时的 \(17.9^{\circ}\) 到远日点时的 \(27.8^{\circ}\) 不等\(^\text{[143][146]}\)。在最大西伸角时，水星在日出前达到最早的升起时间；而在最大东伸角时，水星在日落后达到最晚的落下时间\(^\text{[147]}\)。
\begin{figure}[ht]
\centering
\includegraphics[width=8cm]{./figures/8e02a3e95f06ba7b.png}
\caption{卡内基断崖的伪彩色图像，这是一处构造地形——高地（红色）；低地（蓝色）。} \label{fig_Mercur_17}
\end{figure}
与北半球相比，在南半球更常且更容易观测到水星。这是因为水星的最大西伸角只会出现在南半球的早秋，而其最大东伸角则只会出现在南半球的晚冬\(^\text{[147]}\)。在这两种情况下，行星轨道与地平线的夹角达到最大，使得在前一种情况中，水星可在日出前数小时升起；而在后一种情况中，水星可在日落后数小时才落下。这种现象在南半球中纬度地区（如阿根廷与南非）尤为明显\(^\text{[147]}\)。

另一种观测水星的方法是在白天使用望远镜观测行星，只要天气晴朗，最好在其达到最大伸角时进行。这种方法可以轻松找到水星，即便使用口径仅为 \(8\,\mathrm{cm}\)（\(3.1\,\mathrm{in}\)）的望远镜。然而，必须极其小心，务必遮挡太阳以避免严重的眼部伤害\(^\text{[148]}\)。该方法避开了在黄道角度较低（例如秋季傍晚）时只能在薄明时分观测的限制。此时水星位于更高的天空位置，大气扰动对观测影响较小。在接近上合日时，当水星几乎处于最亮状态，它甚至可以在距离太阳仅 \(4^{\circ}\) 的角度被观测到。

水星也可以像其他行星及最亮的恒星一样，在一次日全食期间被观测到\(^\text{[149]}\)。
\subsection{观测历史}
\subsubsection{古代天文学家}
\begin{figure}[ht]
\centering
\includegraphics[width=8cm]{./figures/84dbbb719eca0d35.png}
\caption{《天文学书》（Liber astronomiae，1550 年）中的水星图像} \label{fig_Mercur_18}
\end{figure}
目前已知最早关于水星的记录来自 MUL.APIN 泥版。这些观测最可能由一位亚述天文学家于公元前 14 世纪左右完成\(^\text{[150]}\)。MUL.APIN 泥版中用于指称水星的楔形文字名称被转写为 UDU.IDIM.GU\U4.UD，意为“跳跃的行星”\(^\text{[c][151]}\)。巴比伦关于水星的记录可追溯至公元前 1 千年，巴比伦人将该行星称为 Nabu，以他们神话中诸神使者之名命名\(^\text{[152]}\)。

希腊—埃及天文学家托勒密（Ptolemy）在其著作《行星假说》（Planetary Hypotheses）中讨论了行星从太阳表面经过（凌日）的可能性。他提出之所以没有观测到这种现象，要么是因为像水星这样的行星过于渺小而无法看见，要么是因为凌日事件的发生过于罕见\(^\text{[153][154]}\)。
\begin{figure}[ht]
\centering
\includegraphics[width=8cm]{./figures/d58620333a1f9d4d.png}
\caption{伊本·沙提尔（Ibn al-Shatir）用于解释水星观测现象的模型，展示了通过图西对（Tusi couple）实现本轮的倍增，从而消除托勒密体系中的偏心圆与均差点。} \label{fig_Mercur_19}
\end{figure}
在中国古代，水星被称为“辰星”，与五行体系中“水”之相及北方方位相关\(^\text{[155]}\)。现代汉语、韩语、日语与越南语文化中，均依据五行体系将其字面称为“水星”\(^\text{[156][157][158]}\)。在印度神话中，水星被称为 Budha，这位神被认为主管星期三\,[159]。日耳曼异教中的奥丁（Odin，或 Woden）亦与水星及星期三相关\(^\text{[160]}\)。玛雅文明可能以猫头鹰（或可能是四只猫头鹰——两只代表晨星相，两只代表昏星相）来象征水星，猫头鹰被视为通往冥界的使者\(^\text{[161]}\)。水星有时也被称为 Stilbon（希腊语：Στίλβων），意为“闪耀者、发光者”\(^\text{[162]}\)。

在中世纪伊斯兰天文学中，11 世纪安达卢西亚天文学家阿布·伊沙克·易卜拉欣·扎尔卡里（Abū Ishāq Ibrāhīm al-Zarqālī）描述了水星地心轨道的均轮为椭圆形，如同鸡蛋或松子一般，尽管这一洞见并未影响他的天文理论或计算\(^\text{[163][164]}\)。12 世纪，伊本·巴贾（Ibn Bajjah）曾观测到“两个黑点出现在太阳表面”，13 世纪马拉盖天文学家库特卜·丁·希拉兹（Qotb al-Din Shirazi）认为这可能是水星和/或金星的凌日现象\(^\text{[165]}\)。然而，中世纪此类报告后来多被解释为对太阳黑子的观测\(^\text{[166]}\)。

在印度，15 世纪喀拉拉学派的天文学家尼拉坎塔·索马亚吉（Nilakantha Somayaji）提出了一种部分日心的行星模型，其中水星绕太阳运行，而太阳再绕地球运行，与 16 世纪晚期第谷·布拉赫（Tycho Brahe）后来的“第谷体系”相似\(^\text{[167]}\)。

1965 年，意大利天文学家朱塞佩·科伦坡（Giuseppe Colombo）注意到水星的自转周期约为其公转周期的三分之二，并提出水星的公转与自转周期不是以 \(1:1\) 的方式潮汐锁定，而是陷入了 \(3:2\) 共振\(^\text{[180]}\)。随后水手 10 号（Mariner~10）的数据证实了这一观点\(^\text{[181]}\)。这意味着夏帕瑞利（Schiaparelli）和安东尼亚迪（Antoniadi）的地图并非“错误”。实际上，天文学家在每隔一个轨道周期时会看到相同的地貌特征并将其记录下来，而在中间那个轨道（当水星另一面对着太阳）进行的观测则因轨道几何导致的极差观测条件而被忽略\(^\text{[171]}\)。

尽管地基光学观测并未显著推进对水星的理解，但使用微波波段干涉测量的射电天文学家——该技术可消除太阳辐射影响——能够探测水星数米深的地下层的物理与化学性质\(^\text{[182][183]}\)。直到第一艘空间探测器飞掠水星后，人们才真正得知该行星许多最基本的形貌特征。此外，技术进步也改善了地基观测能力。2000 年，威尔逊山天文台（Mount Wilson Observatory）使用 1.5 米（4.9 英尺）Hale 望远镜进行了高分辨率幸运成像（lucky imaging）观测，首次解析了水星表面那些未被水手 10 号拍摄区域的地形特征\(^\text{[184]}\)。目前，水星的大部分区域已由阿雷西博射电望远镜以 5 公里（3.1 英里）的分辨率完成测绘，包括可能含有水冰的永久阴影陨石坑中的极区沉积物\(^\text{[185]}\)。











